% Generic one-page resume template.
%
% Build with (from repo root):
%   python3 scripts/build.py --template
% or directly:
%   tectonic templates/resume.tex
%
% OPTIONAL PHOTO (top-right):
%   Drop a headshot at assets/photo.jpg (.jpeg/.png also work), then build
%   with scripts/build.py. If no photo is found, the header automatically
%   falls back to a centered layout with no empty box.
%   To move the photo to the top-LEFT instead, set \photolefttrue below.

\documentclass[11pt, a4paper]{article}
\usepackage[a4paper, top=1.6cm, bottom=1.6cm, left=1.7cm, right=1.7cm]{geometry}
\usepackage{fontspec}
\usepackage[english]{babel}
\usepackage{enumitem}
\usepackage{titlesec}
\usepackage{graphicx}
\usepackage[hidelinks]{hyperref}

\setlist[itemize]{label=-, leftmargin=*, topsep=2pt, partopsep=0pt, parsep=0pt, itemsep=2pt}
\titleformat{\section}{\large\bfseries\uppercase}{}{0em}{}[\titlerule]
\titlespacing*{\section}{0pt}{1.3ex plus 1ex minus .2ex}{0.8ex plus .2ex}

% --- Photo auto-detection ---
% Priority: assets/photo.jpg first. Checked last = wins, so the list below
% runs from lowest to highest priority.
\newcommand{\resumephoto}{}
\newif\ifhasphoto
\hasphotofalse
\newcommand{\tryphoto}[1]{\IfFileExists{#1}{\renewcommand{\resumephoto}{#1}\hasphototrue}{}}
\tryphoto{photo.png}
\tryphoto{photo.jpeg}
\tryphoto{photo.jpg}
\tryphoto{../assets/photo.png}
\tryphoto{../assets/photo.jpeg}
\tryphoto{../assets/photo.jpg}
\tryphoto{assets/photo.png}
\tryphoto{assets/photo.jpeg}
\tryphoto{assets/photo.jpg}

% Flip this to \photolefttrue to put the photo top-left instead of top-right.
\newif\ifphotoleft
\photoleftfalse

\begin{document}
\pagestyle{empty}

% --- HEADER (photo-aware) ---
\ifhasphoto
  % Photo found: contact block (wide) + headshot (narrow, top-aligned).
  \noindent
  \ifphotoleft
    \begin{minipage}[c]{0.20\textwidth}
      \includegraphics[width=\linewidth, height=3.4cm, keepaspectratio, clip]{\resumephoto}
    \end{minipage}\hfill
    \begin{minipage}[c]{0.76\textwidth}
  \else
    \begin{minipage}[c]{0.76\textwidth}
  \fi
      \begin{center}
          {\Huge \textbf{YOUR NAME}} \\[0.2cm]
          {\large Target Role 1 \enskip $\cdot$ \enskip Target Role 2 \enskip $\cdot$ \enskip Target Role 3} \\[0.15cm]
          your.email@example.com \enskip $\cdot$ \enskip +1 000 000 0000 \enskip $\cdot$ \enskip City, Country \\[0.1cm]
          \href{https://github.com/yourusername}{GitHub} \enskip $\cdot$ \enskip
          \href{https://linkedin.com/in/yourusername}{LinkedIn} \enskip $\cdot$ \enskip
          \href{https://your-portfolio.example.com}{Portfolio}
      \end{center}
    \end{minipage}
  \ifphotoleft\else\hfill
    \begin{minipage}[c]{0.20\textwidth}
      \includegraphics[width=\linewidth, height=3.4cm, keepaspectratio, clip]{\resumephoto}
    \end{minipage}
  \fi
  \vspace{0.15cm}
\else
  % No photo found: clean centered header, no empty box.
  \begin{center}
      {\Huge \textbf{YOUR NAME}} \\[0.2cm]
      {\large Target Role 1 \enskip $\cdot$ \enskip Target Role 2 \enskip $\cdot$ \enskip Target Role 3} \\[0.15cm]
      your.email@example.com \enskip $\cdot$ \enskip +1 000 000 0000 \enskip $\cdot$ \enskip City, Country \\[0.1cm]
      \href{https://github.com/yourusername}{GitHub} \enskip $\cdot$ \enskip
      \href{https://linkedin.com/in/yourusername}{LinkedIn} \enskip $\cdot$ \enskip
      \href{https://your-portfolio.example.com}{Portfolio}
  \end{center}
\fi

% --- EDUCATION ---
\section*{Education}
\noindent
\textbf{University Name} \hfill City, Country \\
Degree or Program \hfill Start Year -- End Year

% --- TECHNICAL SKILLS ---
\section*{Technical Skills}
\begin{itemize}
    \item \textbf{Languages:} Python, TypeScript, JavaScript, C++, SQL
    \item \textbf{Frameworks \& Tools:} React, Node.js, Docker, Git, Linux
    \item \textbf{Core Areas:} Full-stack Development, Machine Learning, Systems, Cloud
\end{itemize}

% --- EXPERIENCE ---
\section*{Experience}
\noindent
\textbf{Company or Organization} \hfill City, Country \\
\textit{Role Title} \hfill Month Year -- Month Year
\begin{itemize}
    \item Built or improved a system, product, or workflow; include scale, users, latency, cost, or accuracy when possible.
    \item Collaborated with a team or stakeholder group to deliver a measurable technical outcome.
\end{itemize}

\noindent
\textbf{Company or Organization} \hfill City, Country \\
\textit{Role Title} \hfill Month Year -- Month Year
\begin{itemize}
    \item Describe the problem, your action, and the result in one concise line.
    \item Start bullets with strong verbs and avoid listing responsibilities without impact.
\end{itemize}

% --- PROJECTS ---
\section*{Projects}
\noindent
\textbf{Project Name} $|$ \textit{Tech Stack}
\begin{itemize}
    \item Explain what the project does, the technical choices you made, and why they mattered.
    \item Add a metric, deployment detail, benchmark, user count, or concrete result when available.
\end{itemize}

\noindent
\textbf{Project Name} $|$ \textit{Tech Stack}
\begin{itemize}
    \item Focus on engineering depth: architecture, performance, security, reliability, or product quality.
    \item Link to public demos or repositories only when they are polished enough to support your application.
\end{itemize}

\vspace{0.05cm}
\noindent
\textit{\small More projects available on \href{https://github.com/yourusername}{GitHub}.}

% --- CERTIFICATIONS ---
\section*{Certifications}
\begin{itemize}
    \item \href{https://example.com/certificate}{Certification Name} -- Issuing Organization
    \item Certification Name -- Issuing Organization
\end{itemize}

% --- LEADERSHIP OR ACTIVITIES ---
\section*{Leadership}
\noindent
\textbf{Organization} \\
\textit{Role} \hfill Month Year -- Month Year
\begin{itemize}
    \item Led, organized, mentored, or shipped something meaningful; keep the language specific and verifiable.
\end{itemize}

\end{document}
